\originalchapter{化学势与巨正则系综}\label{ch:grand}
\originalsection{粒子交换与化学平衡}
温度决定能量交换的平衡, 化学势决定粒子交换的平衡. 定义$\mu=(\partial U/\partial N)_{S,V}$, 同时有$\mu=(\partial F/\partial N)_{T,V}$和$\mu=(\partial G/\partial N)_{T,P}$. 固定变量不同, 这几个导数表示的状态函数也不同. 用$U_N|_{T,V}$代替定义一般会错误.

两个系统交换少量能量与粒子, 合成孤立体系的熵变为
\[
\dd S_{\rm tot}=\left(\frac1{T_2}-\frac1{T_1}\right)\dd U_2
+\left(\frac{\mu_1}{T_1}-\frac{\mu_2}{T_2}\right)\dd N_2.
\]
允许两种独立交换时, 平衡要求$T_1=T_2$及$\mu_1=\mu_2$. 等温时粒子自发由高化学势流向低化学势. 多组分反应若化学计量数为$\nu_i$, 平衡条件是$\sum_i\nu_i\mu_i=0$. 带电粒子还应把电势能计入电化学势.

均匀短程体系的$U(S,V,N)$广延, Euler关系给$U=TS-PV+\mu N$, 因而$G=\mu N$. 对它求微分并与第一律比较, 得Gibbs--Duhem关系$S\dd T-V\dd P+N\dd\mu=0$. 在等温共存两相间有$\dd\mu=v\dd P$, $v=V/N$. 这将用于相变构造.

\begin{example}
稀薄三维气体与表面二维吸附气体平衡. 表面每原子结合能为$-\epsilon_0$, 忽略吸附原子相互作用与最大覆盖限制. 求表面数密度$\sigma$与体数密度$n$的关系.
\end{example}
\begin{solution}
三维$z=V/\lambda_T^3$给$\mu_b=\kb T\ln(n\lambda_T^3)$; 二维$z=A\e^{\beta\epsilon_0}/\lambda_T^2$给$\mu_s=-\epsilon_0+\kb T\ln(\sigma\lambda_T^2)$. 两者相等得到$\sigma=n\lambda_T\e^{\beta\epsilon_0}$. 低温吸附增强, 但增大到接近单层覆盖时本模型不再适用, 需加入位点排斥或相互作用.
\end{solution}

\originalsection{巨配分函数与热力学势}
热库兼作粒子库, 展开其熵$S_b(E_{\rm tot}-E,N_{\rm tot}-N)=\text{常数}-E/T+\mu N/T+\cdots$, 得微观态概率
\[
p_{i,N}=\frac{\e^{-\beta(E_{i,N}-\mu N)}}\Xi,\qquad
\Xi(T,V,\mu)=\sum_{N=0}^{\infty}\e^{\beta\mu N}Z_N.
\]
$N=0$真空项必须计入. 巨热力学势$\mathcal G=U-TS-\mu\avg N=-\kb T\ln\Xi$满足
\[
\dd\mathcal G=-S\dd T-P\dd V-\avg N\dd\mu,
\quad\avg N=\kb T\partial_\mu\ln\Xi,
\quad\Var N=\kb T(\partial_\mu\avg N)_{T,V}.
\]
均匀广延体系有$\mathcal G=-PV$, 有表面效应、陷阱或长程引力时不能直接沿用. 在固定$\mu$求$\beta$导数时, $-\partial_\beta\ln\Xi=\avg{E-\mu N}$, 不能遗漏$\mu N$项.

\originalsection{Bose、Fermi与经典极限}
不固定总数后, 独立模式可以分离. 单粒子模式$\alpha$的Fermi占据数为0或1, 因而$\xi_\alpha=1+\e^{-\beta(\epsilon_\alpha-\mu)}$; Bose占据数为任意非负整数, 因而$\xi_\alpha=(1-\e^{-\beta(\epsilon_\alpha-\mu)})^{-1}$, 要求$\mu<\epsilon_0$. 对数求导得
\[
\avg{n_\alpha}=\frac1{\e^{\beta(\epsilon_\alpha-\mu)}\pm1},\quad
\Var n_\alpha=\avg{n_\alpha}(1\mp\avg{n_\alpha}),
\]
上号为Fermi, 下号为Bose. Fermi涨落被排斥压低, Bose涨落被同态聚集提高. 只有$\e^{\beta(\epsilon-\mu)}\gg1$时, 两者才共同退为$\e^{-\beta(\epsilon-\mu)}$. 这是占据稀薄的Maxwell--Boltzmann极限, 不是简单地删去粒子同一性.

给定平均数时$\sum_\alpha\avg{n_\alpha}=N$隐式决定$\mu(T,V,N)$. 能量$U=\sum_\alpha\epsilon_\alpha\avg{n_\alpha}$, 熵可由巨势或直接状态熵得到. Fermion的模式熵是$-\kb[f\ln f+(1-f)\ln(1-f)]$, Boson是$\kb[(1+n)\ln(1+n)-n\ln n]$. 它们的零占据极限都连续, 使用$0\ln0=0$.
